% A.67 candidate: notes only, pending independent full audits.
\subsection{The tetrahedral geometric closures}

\begin{statement}{A.67}{Proved}{All intermediates of every tetrahedral closure}
Every geometrically connected finite cover of
\(U=E\setminus\{O\}\), \(E:y^2=x^3-x\), whose geometric Galois
closure is a quotient of \(A_4\), satisfies both assertions of
Task A.5 at every coefficient prime, every original cusp packet,
and every permitted open derived subgroup.

For an actual \(A_4\)-Galois cover \(X^\circ\to U\), the proper
curve \(X\) has genus four. Its intermediate covers have the
following invariants; subgroups in the same row are conjugate.
\[
\begin{array}{c|r|r|r}
 S\le A_4&[A_4:S]&g(X/S)&\#((X/S)\setminus(X^\circ/S))\\ \hline
 1&12&4&6\\
 C_2&6&2&4\\
 C_3&4&2&2\\
 V_4&3&1&3\\
 A_4&1&1&1
\end{array}
\]
The proof identifies every such closure with an actual cover
constructed from a degree-three isogeny. At residue characteristic
three its pointed model has a rational component containing all
its cusps. The same property is proved for its genus-two
intermediates, including the wild \(C_3\) quotient. This is a
non-nilpotent family; the universal Task A.5 remains open.
\end{statement}

\begin{proof}
We may make a finite number-field extension to define all covers,
deck transformations, isogenies and marked points. At the end we
use the packet and open-derived-subgroup comparison in A.41--A.42,
so this preliminary extension does not restrict the asserted
original packets. We first identify the actual curves and maps.

\paragraph{The isogeny and the functions defining the cover.}
Let \(\phi:E'\to E\) be a degree-three isogeny, with kernel
\(T=\{O,P,-P\}\). A plane embedding by \(|3O|\), followed by
translation of coordinates putting \(P=(0,0)\) and its tangent
at \(y=0\), gives
\[
 E':\quad y^2+a\,xy+b\,y=x^3,\qquad P=(0,0),
 \qquad b(a^3-27b)\ne0.                                      \tag{A.67.1}
\]
Indeed, a point of order three is an inflection point. In the
general Weierstrass equation, passing through \((0,0)\) first
kills \(a_6\), the tangent condition kills \(a_4\), and triple
intersection with \(y=0\) kills \(a_2\). Smoothness at \(P\)
forces \(a_3=b\ne0\). Thus (A.67.1) is a coordinate construction,
not an assumed classification formula.

For a variable point \(Q=(x,y)\), the chord-and-tangent law gives
\[
\begin{split}
 f_0(Q)&=x,&
 f_1(Q)=x(Q+P)&=-\frac{by}{x^2},\\
 f_2(Q)=x(Q-P)&=\frac{b(y+b+ax)}{x^2},&
 f_0f_1f_2&=-b^2.
\end{split}                                                   \tag{A.67.2}
\]
Their divisors are respectively
\[
 P+(-P)-2O,\qquad O+P-2(-P),\qquad O+(-P)-2P.
                                                               \tag{A.67.3}
\]
The two independent odd-divisor vectors in (A.67.3) prove that
adjoining two of the three square roots gives a connected
\(V_4\)-cover. Define \(X\) as the smooth proper normalization
of this function field, choosing roots with
\[
 u_j^2=f_j,\qquad u_0u_1u_2=ib.
                                                               \tag{A.67.4}
\]
Translation by \(P\) cyclically permutes the \(f_j\), so cyclically
permuting the \(u_j\) lifts it to an automorphism \(\tau\) of order
three. The even sign changes of the roots form \(V_4\), and
\[
 V_4\rtimes\langle\tau\rangle=A_4
                                                               \tag{A.67.5}
\]
is the deck group of the actual composite \(X\to E'\to E\).
The only branch points of \(X\to E'\) are \(T\); at each the
inertia has order two, and the three inertia subgroups are
distinct. Thus \(X\to E\) is unramified away from \(O\).
Riemann--Hurwitz gives \(2g(X)-2=3\cdot2=6\), and there are six
points over \(O\). Every nontrivial generic stabilizer is a
subgroup of \(V_4\) of order two.

\paragraph{Why this includes every \(A_4\) closure.}
The cusp of the rank-two geometric fundamental group is a
commutator. In a surjection to \(A_4\) its image is a nonzero
element of \(V_4\): otherwise the two generating images would
commute. The \(C_3=A_4/V_4\) quotient is abelian, kills the cusp,
and therefore extends to a degree-three elliptic isogeny
\(E'\to E\), after choosing an origin. Over \(E'\) its three
cusp inertia images are the three distinct nonzero elements
of \(V_4\), cycled by translation by \(P\).

Here is uniqueness of the equivariant quadratic cover, including
possible unramified twists. Over the geometric field the exact
residue sequence is
\[
 0\longrightarrow H^1(E',\mathbf F_2)
 \longrightarrow H^1(E'\setminus T,\mathbf F_2)
 \longrightarrow \mathbf F_2[T]^0\longrightarrow0.
                                                               \tag{A.67.6}
\]
Translation acts trivially on the first term, of dimension two.
The last term, the sum-zero subspace on the three points, is the
irreducible two-dimensional \(C_3\)-module over \(\mathbf F_2\).
Since \(3\) is invertible in \(\mathbf F_2\), the sequence splits
as \(C_3\)-modules. Consequently its unique nontrivial
two-dimensional isotypic summand is the image of the dual of
any equivariant \(V_4\)-quotient with the specified inertia.
Changing its basis only changes the labels of the deck group.
In particular the trivial unramified summand cannot supply an
additional equivariant quadratic twist. The lift of \(C_3\)
is unique up to \(V_4\)-conjugacy, since
\(H^1(C_3,V_4)=0\), by averaging. Equations (A.67.2)--(A.67.5)
therefore give every geometric \(A_4\) closure, as an actual
cover of \(E\), after a finite field extension.

\paragraph{An exact quotient equation, rather than an assumed \(j\)-formula.}
Summing the translated coordinates and subtracting the kernel
coordinates gives the invariant functions
\[
\begin{split}
 X_0&=x+\frac{ab}{x}+\frac{b^2}{x^2},\\
 Y_0&=\frac{(x^3-abx-2b^2)y-b(ax+b)^2}{x^3}.
\end{split}                                                   \tag{A.67.7}
\]
For example \(X_0=x(Q)+x(Q+P)+x(Q-P)\), since the two finite
kernel \(x\)-coordinates are zero. The addition law gives
\[
 y(Q+P)=\frac{by(y+ax)}{x^3}-b,\qquad
 y(Q-P)=-\frac{b(y+b+ax)^2}{x^3},
\]
and \(Y_0=y(Q)+y(Q+P)+y(Q-P)+b\), proving its invariance too.
Substitution in (A.67.1), without division by a numerical
specialization, verifies the equation
\[
 Y_0^2+aX_0Y_0+bY_0
       =X_0^3-5abX_0-(a^3b+7b^2).
                                                               \tag{A.67.8}
\]
The function \(X_0\) has a double pole at each of the three
kernel points and no others. Its degree is six. The smooth
elliptic equation (A.67.8) has \(X_0\)-degree two, so the map
has degree three and is exactly the quotient by \(T\).

The usual Weierstrass invariant polynomials applied to
(A.67.8) give
\[
 c_4=a(a^3+216b),\qquad
 \Delta=b(a^3-27b)^3.
                                                               \tag{A.67.9}
\]
Set \(r=a^3/b\). The quotient is isomorphic over a finite
extension to the prescribed \(E\), if and only if
\[
 \frac{r(r+216)^3}{(r-27)^3}=1728.
\]
The exact polynomial identity
\[
 r(r+216)^3-1728(r-27)^3=(r^2-540r-5832)^2
                                                               \tag{A.67.10}
\]
therefore yields
\[
                         r=270\pm162\sqrt3.
                                                               \tag{A.67.11}
\]
In particular \(a\ne0\). These computations use the actual
isogeny, not merely a remembered modular correspondence.

\paragraph{The complementary elliptic factor and residue two.}
One of the three double quotients of \(X\to E'\) is obtained by
putting \(z^2=x\) and \(Y=2y+az^2+b\):
\[
 D_0:\quad Y^2=4z^6+a^2z^4+2abz^2+b^2.
                                                               \tag{A.67.12}
\]
Its involution \((z,Y)\mapsto(-z,Y)\) has quotient \(E'\).
The involution \((z,Y)\mapsto(-z,-Y)\) has the actual quotient
\[
 t=z^{-2},\quad w=Y/z^3,\qquad
 F:\ w^2=4+a^2t+2ab t^2+b^2t^3.
                                                               \tag{A.67.13}
\]
After \(X_F=b^2t,\ Y_F=b^2w\), this is
\[
 Y_F^2=X_F^3+2abX_F^2+a^2b^2X_F+4b^4.
                                                               \tag{A.67.14}
\]
The two elliptic differentials pull back to independent
multiples of \(z\,dz/Y\) and \(dz/Y\). Hence
\(J_{D_0}\sim E'\times F\). The three nontrivial \(V_4\)-character
spaces in \(H^1(X)\) are cyclically permuted by the actual
automorphism \(\tau\); each is the complement to \(H^1(E')\)
in one of these double quotients. Consequently
\[
                         J_X\sim E'\times F^3.
                                                               \tag{A.67.15}
\]
This identifies the factor \(F\) using actual maps and
differentials.

For (A.67.14), direct invariant calculation gives
\[
 c_4(F)=16a^2b^2,\quad
 \Delta(F)=256b^7(a^3-27b),\quad
 j(F)=\frac{16r^2}{r-27}=4608\pm2688\sqrt3.
                                                               \tag{A.67.16}
\]
At every place above two, \(j(F)=128(36\pm21\sqrt3)\) has
positive valuation, since the factor in parentheses is a unit.
Thus \(F\) has potentially good reduction and its good special
fibre has \(j=0\), by Proposition VII.5.5, including its
characteristic-two clause in Appendix A.1.4(b)
\cite[pp.~197, 199]{fghSilvermanAEC}. An elliptic curve with
\(j=0\) in characteristic two is supersingular: over the
algebraic closure it is \(y^2+y=x^3\), whose inversion
\((x,y)\mapsto(x,y+1)\) has no affine fixed point. Hence it has
no nonzero geometric point of order two.

The curve \(E'\), being isogenous to \(E\), also has potentially
good supersingular reduction at two. Good reduction is
preserved under isogeny and quotient
\cite[Theorem 1 and Corollary 2, p.~493]{fghSerreTate1968}.
The isogeny (A.67.15) therefore proves that \(J_X\) has potentially
good reduction of two-rank zero. Every intermediate Jacobian
is an actual quotient of \(J_X\), by the norm map, and has the
same two properties. Statement A.46 supplies the pro-\(2\)
normalizer assertion for each of these curves.

\paragraph{A three-adic marked base model.}
At a place above three, make a finite extension so that \(E'\)
has good supersingular reduction. This is possible because it
is isogenous to \(E\); supersingularity is preserved by the
specialized isogeny. The three points \(O,P,-P\) all reduce to
\(O\). In a good formal parameter \(t\), put \(c=t(P)\). Then
\[
 v(c)>0,\qquad t(-P)=-c+O(c^2).
                                                               \tag{A.67.17}
\]
In the cluster chart \(t=c\,u\), the three marks have distinct
reductions \(0,1,-1\). The integral formal addition law shows
that translation by \(P\) reduces to \(u\mapsto u+1\):
all terms of total degree at least two become divisible by
\(c\) after division by \(c\). The marked base thus consists of
the good elliptic component and a rational component with
marks \(0,1,-1\), attached at \(u=\infty\); regular resolution
may insert rational chains. No assertion about a common
radius of all three-torsion points is needed.

Each function \(f_j\) has degree-zero horizontal divisor on
the rational component, no horizontal divisor on the main
component, and no other horizontal zeros or poles. The
intersection equations for its principal divisor imply that
its vertical orders are constant along this tree. Explicitly,
the graph Laplacian of the vertical-order vector is zero,
because the horizontal degree on every component is zero;
the kernel of the Laplacian of a connected finite tree consists
of constants. After multiplying by a constant and extending
the field to extract its square root, the vertical orders
are therefore zero. On the main elliptic component each
reduced \(f_j\) is a nonzero constant. On the rational
component their reduced divisors, and hence their functions
up to nonzero constants, are
\[
 \frac{u^2-1}{u^2},\qquad
 \frac{u(u-1)}{(u+1)^2},\qquad
 \frac{u(u+1)}{(u-1)^2}.
                                                               \tag{A.67.18}
\]
All constants may be removed after a finite extension.

Since two is invertible, Hensel's lemma splits the normalization
over the elliptic component into four copies of the good
elliptic curve. At the attaching node all three functions are
units and have zero slopes, so the same quadratic covers are
étale and split there. Over the rational component, (A.67.18)
defines a connected tame \(V_4\)-cover branched at \(0,1,-1\),
with three distinct inertia groups of order two. Its genus is
zero by Riemann--Hurwitz, and there are four unramified points
over infinity. Thus normalization of the actual global
function field, followed by regular resolution, gives
\[
 \text{one rational centre, four supersingular elliptic tails,
 and six distinct cusp sections on the centre.}                \tag{A.67.19}
\]
The only possible inserted components are rational chains.
These local normalizations agree on the node collars because
they are restrictions of the same Kummer field (A.67.4).
The four tail genera already sum to \(g(X)=4\); the graph is
a tree. This also proves potentially good reduction of \(J_X\)
with three-rank zero.

\paragraph{The exact residual centre and its automorphisms.}
The preceding rational cover can be written over the algebraic
closure of \(\mathbf F_3\) as
\[
             u=\frac{(v^2+1)^2}{v^3-v}.
                                                               \tag{A.67.20}
\]
Direct substitution gives
\[
 u(v+1)=u(v)+1,\qquad u(-1/v)=u(v).
                                                               \tag{A.67.21}
\]
The conjugates of \(v\mapsto-1/v\) generate \(V_4\);
\(v\mapsto v+1\) gives the order-three lift. Any other such
lift is conjugate to it by \(V_4\), as already proved by
\(H^1(C_3,V_4)=0\). The four nodes are
\(\mathbf P^1(\mathbf F_3)\). The derivative numerator is
\[
 v^6+v^4+v^2+1=\frac{v^9-v}{v^3-v},
                                                               \tag{A.67.22}
\]
so the six cusps are exactly \(\mathbf F_9\setminus\mathbf F_3\).
They have stabilizer of order two in
\(\operatorname{PSL}_2(\mathbf F_3)=A_4\).

In particular a \(C_3\) fixes one of the four nodes and cycles
the other three. It has no fixed cusp. Its quotient coordinate
\[
                              w=v^3-v
                                                               \tag{A.67.23}
\]
sends the two cusp orbits to \(i,-i\), where \(i^2=-1\), and
the two node orbits to \(0,\infty\). A subgroup \(C_2\le V_4\)
fixes two cusps and no node. Its quotient centre has four
cusp images and two node images. All of these images are
distinct as indicated.

\paragraph{The tame genus-two intermediates.}
For \(S=C_2\), quotienting the actual three-adic model is tame.
The four tails form two free \(C_2\)-orbits, giving two elliptic
tails. The rational centre has the four distinct marks and
two nodes just computed. Local quotient singularities at
nodes are resolved by rational chains. Thus every cusp lies
on this rational marked centre, and joining two cusps does
not pass through a positive-genus component. Its genus is
two, also directly from Riemann--Hurwitz on \(X\to X/C_2\).

\paragraph{The wild genus-two intermediate: residue fields and resolution.}
Put \(D=X/C_3\). The action is free on the generic proper curve,
since the only nontrivial generic stabilizers have order two.
Hence \(g(D)=2\), and its two cusps are the two \(C_3\)-orbits
already computed. The norm \(J_X\to J_D\) is surjective. Since
\(J_X\) has good reduction after the chosen extension, so does
\(J_D\), by the good-reduction criterion cited above. Therefore
any semistable model of \(D\) has a tree as its dual graph.

Take the finite normal quotient of the actual model of \(X\)
by \(C_3\), and then resolve and semistabilize, after a further
finite extension if necessary. We do not assume that forming
this quotient commutes with reduction. Instead, consider the
component valuations and their residue fields. Distinct
orbits of components give distinct valuations of the quotient
field. On the rational centre the residual action is
\(v\mapsto v+1\). The quotient residue field is contained in
\(k(v)^{C_3}=k(v^3-v)\). If \(A\) is an integral lift of \(v\)
at this valuation, the coefficients of
\[
                      \prod_{g\in C_3}(T-gA)
\]
belong to the quotient field and are integral. Their
reductions give \(T^3-T-(v^3-v)\), so \(v^3-v\) belongs to the
quotient residue field. The centre image is consequently
the rational component with coordinate (A.67.23), not an
unspecified reduction of a quotient.

The orbit of three elliptic tails has trivial decomposition
group and gives one elliptic component isomorphic to a tail.
For the fixed tail, the map \(X\to E'\) is an isomorphism on
that component. Translation by \(P\) on the good model of
\(E'\) reduces to the identity because \(P\) reduces to \(O\).
Thus the residual action on the fixed elliptic tail is
trivial. For any residue function \(a\) on that tail, lift it
to an integral element \(A\). The invariant polynomial above
now reduces to \((T-a)^3=T^3-a^3\), proving that the quotient
residue field contains all cubes. It lies between \(k(C)^3\)
and \(k(C)\), where \(C\) is the elliptic tail. Since \(k\)
is perfect and \(k(C)\) has transcendence degree one, this
extension is purely inseparable of degree at most three.
The normalization of the image is therefore an elliptic
curve, either \(C\) or its Frobenius twist.

The two tail orbits have thus produced two distinct
positive-genus components, each of genus one. Their genera
exhaust \(g(D)=2\). The graph is a tree by good reduction of
\(J_D\); hence all further components in resolution and
semistabilization are rational. This genus argument is what
excludes extra positive-genus exceptional components of the
wild quotient.

Finally the two cusp orbits are free already on the residual
centre. Near each such orbit the quotient is the quotient of
three disjoint smooth neighbourhoods permuted freely, so the
marked images remain smooth and distinct. They are at
\(w=i,-i\), whereas the two node orbits are at \(0,\infty\).
The two nodes lead to the two elliptic image components,
possibly through rational chains. Thus the rational centre
has two marks and two distinct outgoing branches and survives
marked stabilization. Both cusps lie on it, and their joining
path is rational. This proves the exact pointed hypothesis
needed for \(D\), including the wild order-three comparison.
The resolution and stabilization operations used here are
the usual point blow-ups and contraction of unstable rational
components; see \cite[9.2.5, 9.2.7, 10.4.24]{fghLiu}.

\paragraph{Fixed subgroups and the remaining normalizer primes.}
For \(X\) and both types of genus-two intermediate, the
three-adic constructions give a compact-type pointed model
in which every cusp has another cusp joined to it entirely
through rational components. Apply A.41--A.42 to conclude
that every permitted fixed subgroup, at every coefficient
prime, is exactly the chosen cusp inertia. The two
genus-one rows are A.40.

For normalizers at coefficient prime three, all intermediate
Jacobians have potentially good reduction of three-rank zero:
this follows from (A.67.19) and the actual norm quotients.
Apply A.46. The coefficient-two assertion was proved using
(A.67.15)--(A.67.16) and A.46. At any coefficient prime
\(\ell\ge5\), the geometric Galois closure has order dividing
twelve, prime to \(\ell\). The good one-marked elliptic model
of \(E\), after a finite extension, admits the actual
prime-to-\(\ell\) tame extension of this cover. The
normalization is a smooth proper curve with the prescribed
marked branch fibres; equivalently this is the specialization
equivalence of SGA~1, Expos\'e XIII, 2.9--2.12
\cite{fghSGA1}, as in A.59. Every intermediate proper curve
therefore has potentially good curve reduction at \(\ell\).
Apply A.57 for genus at least two and A.40 in genus one.
This proves the normalizer assertion for all primes.

The subgroup table follows directly from the subgroup types
of \(A_4\), the six-point cusp action, and Riemann--Hurwitz;
the exact identities \(2g-2+r=\deg(Y/U)\) provide an
additional check. Proper quotients of \(A_4\) are trivial or
\(C_3\), whose covers are elliptic isogenies and are A.40.
The finite-extension comparison of A.41--A.42 now returns
both assertions to every original packet and permitted
open derived subgroup, completing the statement.

The executed exact certificate is
\path{notes/fgh_A67_certificate.py}. It verifies the translation
and three-isogeny identities modulo the actual cubic equation,
the complementary elliptic invariants, both values of \(r\),
the residual rational map and all \(A_4\) orbits on its four
nodes and six cusps. The geometric arguments above, including
the wild quotient and the passage to all covers, are separate
from these finite algebraic checks.
\end{proof}
